% GENERATED FILE. Edit canonical lessons, then run npm run generate:books.
% canonical-source-hash: fnv1a64:ec7ac5098777e54c
% canonical-lessons: ZH-C169-yidazao, ZH-C169-wujiao, ZH-W169-jiao-sleep, ZH-C169-bangwan, ZH-C169-aoye, ZH-C169-shenye, ZH-R169-review-164-169

\chapter{Early and Late}
\label{ch:zh-early-and-late}
\hlchaptermodality{169}

\begin{chapteropening}
\textbf{By the end of this chapter:} \emph{I can talk about first thing in the morning, an afternoon nap, dusk, staying up and the dead of night.}

Complete the last lesson of chapter 169: I can talk about first thing in the morning, an afternoon nap, dusk, staying up and the dead of night.
\end{chapteropening}

\section[\texorpdfstring{\zh{一大早} (yídàzǎo)}{一大早 (yídàzǎo)}]{\zh{一大早} (yídàzǎo) — first thing in the morning; at the crack of dawn}

\label{lesson:ZH-C169-yidazao}



\begin{quote}
Before the new one: say the Mandarin for to boil water, then the Mandarin for sugar.
\end{quote}

\subsection*{You'll want to know: \zh{一大早}}
\textbf{\zh{一大早}} — \emph{yídàzǎo} — \textquotedblleft{}first thing in the morning; at the crack of dawn\textquotedblright{}.

\begin{quote}
\textbf{\zh{一大早妈妈去买大米了}} \emph{yídàzǎo māma qù mǎi dàmǐ le} — First thing this morning my mum went out to buy rice.
\end{quote}

The tone of \zh{一} changes: yī is said yí before the fourth tone of dà.

\begin{cousinweb}
\textbf{\zh{一大早}} (\emph{yídàzǎo}) is built from \zh{一} 'one', \zh{大} 'big' and \zh{早} 'early', all three inherited from Middle Chinese, \emph{'jit}, \emph{daH} and \emph{tsawX} (Guangyun, in Baxter's transcription), which glosses the last 'morning'.
\end{cousinweb}

\begin{grammarlens}[title={What you've built}]
One more word for times of day.
\end{grammarlens}

\subsection*{Your turn}
\begin{itemize}
\raggedright
  \item \emph{Say it:} \emph{yídàzǎo}
  \item \emph{Say it:} \emph{yídàzǎo}, once more
  \item \emph{Say it:} \emph{yídàzǎo māma qù mǎi dàmǐ le} — First thing this morning my mum went out to buy rice.
  \item \emph{Recall:} say the Mandarin for to boil water, then the Mandarin for sugar, then say \emph{yídàzǎo} again
\end{itemize}

\begin{tcolorbox}[breakable,colback=teal!4,colframe=teal!35!black,title={Before you move on}]
What does \zh{一大早} mean? (\textquotedblleft{}First thing in the morning; at the crack of dawn\textquotedblright{}.) Say it once more.
\end{tcolorbox}
\section[\texorpdfstring{\zh{午觉} (wǔjiào)}{午觉 (wǔjiào)}]{\zh{午觉} (wǔjiào) — afternoon nap}

\label{lesson:ZH-C169-wujiao}



\begin{quote}
Before the new one: say the Mandarin for sugar, then the Mandarin for first thing in the morning.
\end{quote}

\subsection*{You'll want to know: \zh{午觉}}
\textbf{\zh{午觉}} — \emph{wǔjiào} — \textquotedblleft{}afternoon nap\textquotedblright{}.

\begin{quote}
\textbf{\zh{午觉以后我们喝茶}} \emph{wǔjiào yǐhòu wǒmen hē chá} — After the nap we'll have tea.
\end{quote}

The new character, in pieces: \zh{觉} — nine strokes: three short strokes on top, the cover piece of \zh{学} under them, and \zh{见} at the foot. The next lesson puts it in your hand.

\begin{cousinweb}
\textbf{\zh{午觉}} (\emph{wǔjiào}) is built from \zh{午} 'noon' and \zh{觉} 'sleep', both inherited from Middle Chinese, \emph{nguX} and \emph{kaewH} (Guangyun, in Baxter's transcription), which glosses the second with the word for sleeping. The same character read \emph{jué}, as in \zh{觉得}, is its other reading, glossed 'to wake' and 'to know'.
\end{cousinweb}

\begin{grammarlens}[title={What you've built}]
One more word for times of day.
\end{grammarlens}

\subsection*{Your turn}
\begin{itemize}
\raggedright
  \item \emph{Say it:} \emph{wǔjiào}
  \item \emph{Say it:} \emph{wǔjiào}, once more
  \item \emph{Say it:} \emph{wǔjiào yǐhòu wǒmen hē chá} — After the nap we'll have tea.
  \item \emph{Recall:} say the Mandarin for sugar, then the Mandarin for first thing in the morning, then say \emph{wǔjiào} again
\end{itemize}

\begin{tcolorbox}[breakable,colback=teal!4,colframe=teal!35!black,title={Before you move on}]
What does \zh{午觉} mean? (\textquotedblleft{}Afternoon nap\textquotedblright{}.) Say it once more.
\end{tcolorbox}
\section[\texorpdfstring{\zh{觉} (jiào)}{觉 (jiào)}]{\zh{觉} — three strokes, a cover, then \zh{见}}

\label{lesson:ZH-W169-jiao-sleep}



\begin{quote}
The word first, then the character inside it.

\begin{itemize}
\raggedright
  \item \emph{Recall:} say \emph{wǔjiào} — \textbf{R1}, one lesson back, before it decays
\end{itemize}

\begin{quote}
\zh{见} \zh{学}
\end{quote}

\emph{Read it:} each one aloud, with nothing said first — \emph{jiàn}, \emph{xué}
\end{quote}

\subsection*{Script — three strokes, a cover, then \zh{见}}
\hlblockfigure{\detokenize{figures/ZH-W169-jiao-sleep-filmstrip.pdf}}{How \zh{觉} is written, stroke by stroke}

\begin{quote}
\zh{觉}
\end{quote}

\textbf{Three pieces, top to bottom.} Count the pieces, not the strokes.

\textbf{Three short strokes on top:} a dot on the left, a dot in the middle, and a short stroke on the right falling to the left.

\textbf{Then the cover piece of \zh{学}, two strokes:} a short stroke falling down on the left, then a horizontal that turns down and hooks.

\textbf{Then \zh{见}, four strokes,} as you write it, under the cover.

3 + 2 + 4 = nine.

\subsection*{Your turn}
\begin{itemize}
\raggedright
  \item \emph{Write it:} \zh{觉}, all nine, with the model in front of you
  \item \emph{Cover:} write it again from memory, and count the pieces, not the strokes
\end{itemize}

\begin{tcolorbox}[breakable,colback=teal!4,colframe=teal!35!black,title={Before you move on}]
How many strokes in \zh{觉}? (\textbf{Nine}.) What sits under the cover? (\textbf{\zh{见}}.) Which word holds the cover too? (\textbf{\zh{学}}.)

Source: \href{https://raw.githubusercontent.com/chanind/hanzi-writer-data/68d10a4b21150cae5e1ebbd223eed289cf32d90c/data/\%E8\%A7\%89.json}{Hanzi Writer Data \zh{觉}, pinned stroke paths and medians}.
\end{tcolorbox}
\section[bàngwǎn]{bàngwǎn — dusk; towards evening}

\label{lesson:ZH-C169-bangwan}



\begin{quote}
Before the new one: say the Mandarin for first thing in the morning, then the Mandarin for an afternoon nap.
\end{quote}

\subsection*{You'll want to know: bàngwǎn}
\emph{bàngwǎn} — \textquotedblleft{}dusk; towards evening\textquotedblright{}.

\begin{quote}
\emph{wǒmen bàngwǎn qù sànbù} — We go for a walk around dusk.
\end{quote}

Its characters come later; for now it is a word for the ear.

\begin{cousinweb}
\emph{bàngwǎn} is built from \emph{bàng} 'near' and \zh{晚} 'evening', both inherited from Middle Chinese, \emph{bangH} and \emph{mjonX} (Guangyun, in Baxter's transcription). The dictionary gives the sense 'near' under the level-tone reading of the first; today's falling tone goes with \emph{bangH}, which it leaves without a gloss.
\end{cousinweb}

\begin{grammarlens}[title={What you've built}]
One more word for times of day.
\end{grammarlens}

\subsection*{Your turn}
\begin{itemize}
\raggedright
  \item \emph{Say it:} \emph{bàngwǎn}
  \item \emph{Say it:} \emph{bàngwǎn}, once more
  \item \emph{Say it:} \emph{wǒmen bàngwǎn qù sànbù} — We go for a walk around dusk.
  \item \emph{Recall:} say the Mandarin for first thing in the morning, then the Mandarin for an afternoon nap, then say \emph{bàngwǎn} again
\end{itemize}

\begin{tcolorbox}[breakable,colback=teal!4,colframe=teal!35!black,title={Before you move on}]
What does bàngwǎn mean? (\textquotedblleft{}Dusk; towards evening\textquotedblright{}.) Say it once more.
\end{tcolorbox}
\section[áoyè]{áoyè — to stay up late; to stay up all night}

\label{lesson:ZH-C169-aoye}



\begin{quote}
Before the new one: say the Mandarin for an afternoon nap, then the Mandarin for towards evening.
\end{quote}

\subsection*{You'll want to know: áoyè}
\emph{áoyè} — \textquotedblleft{}to stay up late; to stay up all night\textquotedblright{}.

\begin{quote}
\emph{wǒ zuótiān áoyè le jīntiān hěn lèi} — I stayed up late last night, and I'm tired today.
\end{quote}

Its characters come later; for now it is a word for the ear.

\begin{cousinweb}
\emph{áoyè} is built from \emph{áo} 'to simmer' and \emph{yè} 'night', both inherited from Middle Chinese, \emph{ngaw} and \emph{yaeH} (Guangyun, in Baxter's transcription); the dictionary glosses the first 'to simmer'. How the pair came to mean staying up is not settled.
\end{cousinweb}

\begin{grammarlens}[title={What you've built}]
One more word for times of day.
\end{grammarlens}

\subsection*{Your turn}
\begin{itemize}
\raggedright
  \item \emph{Say it:} \emph{áoyè}
  \item \emph{Say it:} \emph{áoyè}, once more
  \item \emph{Say it:} \emph{wǒ zuótiān áoyè le jīntiān hěn lèi} — I stayed up late last night, and I'm tired today.
  \item \emph{Recall:} say the Mandarin for an afternoon nap, then the Mandarin for towards evening, then say \emph{áoyè} again
\end{itemize}

\begin{tcolorbox}[breakable,colback=teal!4,colframe=teal!35!black,title={Before you move on}]
What does áoyè mean? (\textquotedblleft{}To stay up late; to stay up all night\textquotedblright{}.) Say it once more.
\end{tcolorbox}
\section[shēnyè]{shēnyè — late at night; the dead of night}

\label{lesson:ZH-C169-shenye}



\begin{quote}
Before the new one: say the Mandarin for towards evening, then the Mandarin for to stay up late.
\end{quote}

\subsection*{You'll want to know: shēnyè}
\emph{shēnyè} — \textquotedblleft{}late at night; the dead of night\textquotedblright{}.

\begin{quote}
\emph{shēnyè wàibian hěn ānjìng} — Late at night it's quiet outside.
\end{quote}

Its characters come later; for now it is a word for the ear.

\begin{cousinweb}
\emph{shēnyè} is built from \emph{shēn} 'deep' and \emph{yè} 'night', both inherited from Middle Chinese, \emph{syim} and \emph{yaeH} (Guangyun, in Baxter's transcription): deep in the night. The dictionary gives 'deep' under the falling-tone reading \emph{syimH}; today's level tone goes with \emph{syim}.
\end{cousinweb}

\begin{grammarlens}[title={What you've built}]
One more word for times of day.
\end{grammarlens}

\subsection*{Your turn}
\begin{itemize}
\raggedright
  \item \emph{Say it:} \emph{shēnyè}
  \item \emph{Say it:} \emph{shēnyè}, once more
  \item \emph{Say it:} \emph{shēnyè wàibian hěn ānjìng} — Late at night it's quiet outside.
  \item \emph{Recall:} say the Mandarin for towards evening, then the Mandarin for to stay up late, then say \emph{shēnyè} again
\end{itemize}

\begin{tcolorbox}[breakable,colback=teal!4,colframe=teal!35!black,title={Before you move on}]
What does shēnyè mean? (\textquotedblleft{}Late at night; the dead of night\textquotedblright{}.) Say it once more.
\end{tcolorbox}
\section[(dialogue)]{Looking back: from Getting Ready to Early and Late}

\label{lesson:ZH-R169-review-164-169}



\begin{quote}
Say the Mandarin for to dress up and to wash one's hair.
\end{quote}

\subsection*{Your turn}
\begin{itemize}
\raggedright
  \item \emph{Say it:} to dress up, to wash one's hair, a hair dryer, to put on make-up, to have a shave
  \item \emph{Say it:} an overcoat, a sweater, measure word for long things, a skirt, a scarf
  \item \emph{Say it:} boiled water, afternoon tea, black tea, tea leaves, to make tea
  \item \emph{Say it:} the day before yesterday, to happen, it's nothing, to have a power cut, a sound
  \item \emph{Say it:} to steam, to deep-fry, to stew, to boil water, sugar
  \item \emph{Say it:} first thing in the morning, an afternoon nap, towards evening, to stay up late, late at night
  \item \emph{Read it:} \zh{衣} \zh{茶} \zh{事} \zh{觉} — each one aloud, then a word that holds it
\end{itemize}

\begin{tcolorbox}[breakable,colback=teal!4,colframe=teal!35!black,title={Before you move on}]
To dress up? (\emph{dǎban}.) To wash one's hair? (\emph{xǐtóu}.) A hair dryer? (\emph{chuīfēngjī}.) To put on make-up? (\emph{huàzhuāng}.) To have a shave? (\emph{guāliǎn}.) An overcoat? (\textbf{\zh{大衣}}, \emph{dàyī}.) A sweater? (\emph{máoyī}.) Measure word for long things? (\emph{tiáo}.) A skirt? (\emph{qúnzi}.) A scarf? (\emph{wéijīn}.) Boiled water? (\textbf{\zh{开水}}, \emph{kāishuǐ}.) Afternoon tea? (\textbf{\zh{下午茶}}, \emph{xiàwǔchá}.) Black tea? (\emph{hóngchá}.) Tea leaves? (\emph{cháyè}.) To make tea? (\emph{pàochá}.) The day before yesterday? (\textbf{\zh{前天}}, \emph{qiántiān}.) To happen? (\emph{fāshēng}.) It's nothing? (\textbf{\zh{没事}}, \emph{méishì}.) To have a power cut? (\emph{tíngdiàn}.) A sound? (\emph{shēngyīn}.) To steam? (\emph{zhēng}.) To deep-fry? (\emph{zhá}.) To stew? (\emph{dùn}.) To boil water? (\emph{shāoshuǐ}.) Sugar? (\emph{táng}.) First thing in the morning? (\textbf{\zh{一大早}}, \emph{yídàzǎo}.) An afternoon nap? (\textbf{\zh{午觉}}, \emph{wǔjiào}.) Towards evening? (\emph{bàngwǎn}.) To stay up late? (\emph{áoyè}.) Late at night? (\emph{shēnyè}.)
\end{tcolorbox}
